```latex
\textbf{Proof.}
For each \(r\), the alternating-sign property on
\[
F_r=\{\,t_r+k a_r:0\le k\le m_r\,\}
\]
means that
\[
\sigma(t_r+k a_r)=\sigma(t_r)(-1)^k \qquad (0\le k\le m_r).
\]
Let
\[
t^*:=t_i+p a_i=t_j+q a_j.
\]

\begin{enumerate}
\item Since \(t^*\) is a single element of \(T\), it has a unique sign. Hence
\[
\sigma(t_i)(-1)^p=\sigma(t^*)=\sigma(t_j)(-1)^q.
\]
This is exactly the asserted sign consistency.

\item Rearranging the equality
\[
t_i+p a_i=t_j+q a_j
\]
gives
\[
t_j-t_i=p a_i-q a_j.
\]

\item If \(p\ge 1\) and \(q\ge 1\), then both directions occur nontrivially, and we may write
\[
t_j=t_i+p a_i-q a_j.
\]
Equivalently,
\[
t_i=t^*-p a_i,
\qquad
t_j=t^*-q a_j.
\]
Thus \(t_i\) and \(t_j\) are affinely related through the two edge directions \(a_i\) and \(a_j\).

\item Put
\[
E_r:=[t_r,t_{r+1}]
=\{\,t_r+\lambda a_r:0\le \lambda\le m_r\,\}.
\]
Because
\[
t_{r+1}-t_r=m_r a_r
\]
for every \(r\), and because
\[
\sum_{r=1}^n m_r a_r=0,
\]
the segments \(E_1,\dots,E_n\) form the boundary edges of the convex polygon
\[
Q=\operatorname{conv}(T),
\]
listed in cyclic order.

Also, at the level of signs,
\[
\sigma(t_{r+1})
=\sigma(t_r+m_r a_r)
=\sigma(t_r)(-1)^{m_r},
\]
so along the cyclic boundary arc from \(t_i\) to \(t_j\) one gets
\[
\sigma(t_j)=\sigma(t_i)(-1)^{m_i+\cdots+m_{j-1}}.
\]
Hence a ``sign contradiction'' can only mean a disagreement between this boundary parity and the local parity from part \(1\). Such a disagreement cannot come from a mere shared endpoint, so the only relevant case is the genuine interior-overlap case
\[
0<p<m_i,
\qquad
0<q<m_j.
\]
Then
\[
t^*\in \operatorname{relint}(E_i)\cap \operatorname{relint}(E_j).
\]

After cyclic relabeling, we may assume that \(i<j\). Along the boundary of \(Q\),
\[
t_j-t_i
=\sum_{k=i}^{j-1}(t_{k+1}-t_k)
=\sum_{k=i}^{j-1}m_k a_k.
\]
Combining this with part \(2\), namely
\[
t_j-t_i=p a_i-q a_j,
\]
yields
\[
(m_i-p)a_i+\sum_{k=i+1}^{j-1}m_k a_k+q a_j=0.
\]
Geometrically, this is exactly the vector equation of the boundary subchain
\[
t^*\to t_{i+1}\to t_{i+2}\to \cdots \to t_j\to t^*.
\]

But the boundary of a convex polygon is a simple polygonal curve, and the relative interiors of distinct boundary edges are disjoint. Therefore a proper boundary subchain can close up only if its first and last edge pieces lie on the same supporting line. Consequently \(E_i\) and \(E_j\) must have the same direction. Since both are oriented by the cyclic order of \(\partial Q\), their direction vectors are positive multiples of one another:
\[
a_j=\lambda a_i
\qquad\text{for some }\lambda>0.
\]
Equivalently,
\[
a_i\parallel a_j
\]
with positive proportionality, i.e. \(a_i\) and \(a_j\) lie in the same direction class.

This proves the claim.
\hfill\(\square\)
```
